\chapter*{Introduction générale}
\addcontentsline{toc}{chapter}{Introduction générale}

Le secteur du transport routier de marchandises repose sur une coordination permanente entre clients, chauffeurs, véhicules et missions. Une organisation qui s'appuie uniquement sur des tableurs ou des échanges informels expose l'exploitant à des oublis : mission mal affectée, véhicule indisponible, permis ou assurance arrivés à échéance. Ces ruptures ont un coût opérationnel et réglementaire.

Le projet de stage \textbf{SmartTrans-AI} vise à concevoir et réaliser une application web de gestion de transport, enrichie par un assistant conversationnel. L'objectif n'est pas de remplacer le responsable d'exploitation, mais de centraliser les données, de tracer les actions et d'accélérer les recherches du type : combien de missions sont en cours, quel chauffeur est disponible, quel véhicule peut être affecté.

Le périmètre retenu comprend l'authentification, le tableau de bord, la gestion des utilisateurs (administrateur), des clients, des chauffeurs et des véhicules, le cycle de vie des missions, le journal d'activités, les alertes d'expiration et l'assistant IA. Sont volontairement exclus du produit livré un module documentaire de type GED et un écran de notifications push. Le chauffeur peut consulter ses missions, les démarrer et les terminer ; il ne peut pas les annuler.

Ce rapport s'organise en sept chapitres. Le premier situe le cadre du stage chez S. WATRANS. Le deuxième dresse un état de l'art et positionne le produit. Le troisième présente l'analyse et la conception (cas d'utilisation, classes, données). Le quatrième décrit la réalisation. Le cinquième traite des tests, du déploiement et des limites. Le sixième revient sur la conduite de projet. Le septième présente les interfaces (captures d'écran). Une conclusion générale ouvre des perspectives.
